% !TeX root = main.tex
\lecture{5}{Thu 29 Jan 2026 09:00}{Partial Differentiation V: 3D Stationary Points and Lagrange Multipliers}

\section{2D Recap}
In 2D, a stationary point of a function $f(x)$ occurs at $x = a$ where $f(x)$ is neither increasing or decreasing such that $f^\prime(a) = 0$.

Consider the Taylor series about $x = a$, where $f^\prime(a) = 0$:

\[
    f(a+h) = f(a) + h f^\prime(a) + \frac{h^2}{2!} f^{\prime\prime} (a) + \cdots = f(a) + \frac{h^2}{2!} f^{\prime\prime} (a) + \cdots
\]

If $f^{\prime\prime} (a) > 0$, the function looks like $+h^2$ near the stationary point. This gives us a minima. If the second derivative is negative the function looks like $-h^2$, giving a maximum point.

If the second derivative is zero, we need to find the first non-zero derivative $f^n{(a)} \neq 0$ where $n > 2$.
\begin{itemize}
    \item If $n = 2k + 1$ (odd), the function looks like $\pm h^{2k + 1}$ near the stationary point. This results in a point of inflexion.
    \item If $n = 2k$ (even) and $f^{2k}(a) > 0$, the function looks like $+ h^{2k}$ near the stationary point so that we have a minimum.
    \item If $n = 2k$ (even) and $f^{2k}(a) < 0$, the function looks like $- h^{2k}$ near the stationary point so that we have a maximum.
\end{itemize}

\section{Stationary Points of Several Variables}
\subsection{Two Variables}
Consider a function $f(x, y)$ at a point $(a, b)$. For a stationary point to occur here, the function must be flat in all directions at that point, so all the partial derivatives must be zero:
\[
    f_x(a, b) = f_y(a, b) = 0
\]

To classify it, we again consider the Taylor series:
\[
    f(a+h, b+k) \approx f(a, b) + \frac{1}{2} \left(Ah^2 + 2Bhk + Ck^2\right)
\]
Where $A = f_{xx}(a, b)$, $B = f_{xy}(a,b)$ and $C = f_{yy}(a,b)$. We can then complete the square:

\[
    Ah^2 + Bhk + Ck^2 = A \left(h + \frac{B}{A}k\right)^2 + \left(C - \frac{B^2}{A}\right)k^2
\]

\begin{itemize}
    \item For the stationary point to be a minimum, we need $A > 0$ and $C - B^2 / A = AC - B^2 > 0$.
    \item For a maximum, we need $A < 0$ and $C-B^2/A = AC-B^2 < 0$.
    \item For a saddle point, with a minimum in one direction and a maximum in the other, we need $A$ and $C - B^2/A$ to be oppositely signed. This corresponds to the case where $AC - B^2 < 0$.
\end{itemize}

Considering the second derivatives directly:
\begin{itemize}
    \item For a minimum we need $f_{xx} > 0$ and the Hessian determinant must be positive:
    \[
        \Delta = \det H = \det \begin{bmatrix}
            f_{xx} & f_{xy} \\
            f_{yx} & f_{yy}
        \end{bmatrix} > 0
    \]

    \item For a maximum, we need $f_{xx} < 0$ and the Hessian determinant must be positive:
    \[
        \Delta = \det H = \det \begin{bmatrix}
            f_{xx} & f_{xy} \\
            f_{yx} & f_{yy}
        \end{bmatrix} > 0
    \]
    
    \item For a saddle point, we need the Hessian determinant to be negative:
    \[
        \Delta = \det H = \det \begin{bmatrix}
            f_{xx} & f_{xy} \\
            f_{yx} & f_{yy}
        \end{bmatrix} < 0
    \] 
\end{itemize}

\subsection{Generalising To 3D}
For a function $f(x, y, z)$ to have a stationary point at $(a, b, c)$, all first 
derivatives must be zero. The Taylor expansion then has no linear term, and performing the Taylor Expansion gives:

\[
    f(a+h, b+k, c+l) - f(a,b,c) = \frac{1}{2}\left(Ah^2 + 2Bhk + 2Chl + Dk^2 + 2Ekl + Fl^2\right)
\]

where $A = f_{xx},\, B = f_{xy},\, C = f_{xz},\, D = f_{yy},\, E = f_{yz},\, 
F = f_{zz}$, all evaluated at $(a,b,c)$. We define the quadratic:

\[
    Q = Ah^2 + 2Bhk + 2Chl + Dk^2 + 2Ekl + Fl^2
\]

as the extra factor of $1/2$ does not change the classification of the point. The stationary point is a minimum $Q > 0$ for all $(h, k, l) \neq (0,0,0)$, because a positive $Q$ means that the function increases in all directions away from the point.

Completing the square in $h$ (treating $k$ and $l$ as fixed):
\[
    Ah^2 + 2Bhk + 2Chl = A\!\left(h + \frac{B}{A}k + \frac{C}{A}l\right)^{\!2} 
    - A\!\left(\frac{B}{A}k + \frac{C}{A}l\right)^{\!2}
\]

Expanding the subtracted term:
\[
    -A\!\left(\frac{B}{A}k + \frac{C}{A}l\right)^{\!2} 
    = -A\left(\frac{B^2}{A^2}k^2 + \frac{2BC}{A^2}kl + \frac{C^2}{A^2}l^2\right)
    = -\frac{B^2}{A}k^2 - \frac{2BC}{A}kl - \frac{C^2}{A}l^2
\]

Substituting back and collecting the $k$- and $l$-terms:
\[
    Q = A\!\left(h + \frac{B}{A}k + \frac{C}{A}l\right)^{\!2} 
    + \underbrace{\left(D - \frac{B^2}{A}\right)}_{A'}k^2 
    + 2\underbrace{\left(E - \frac{BC}{A}\right)}_{B'}kl 
    + \underbrace{\left(F - \frac{C^2}{A}\right)}_{C'}l^2
\]

The second quadratic in $A', 2B', C'$ is just our previous 2D example, so that condition holds. For the quadratic, we require:

\[
    A > 0
\]
\[
    D - \frac{B^2}{A} > 0
\]
\[
    \left(D - \frac{B^2}{A}\right) \left(F - \frac{C^2}{A}\right) - \left(E - \frac{BC}{A}\right)^2 > 0
\]

Expanding the last condition we get:
\[
    A^2 DF - ADC^2 - AB^2F + B^2C^2 - E^2 A^2 - B^2 C^2 + 2EABC > 0
\]
\[
    A \left(ADF - DC^2 - B^2 F - E^2 A + 2EBC\right) > 0
\]
\[
    \left(ADF - DC^2 - B^2 F - E^2 A + 2EBC\right) > 0 \qquad \text{ as } A > 0
\]
\[
    A(DF - E^2) - B(BF - EC) + C(BE - CD) > 0
\]
\[
    \implies \det \begin{bmatrix}
        A & B & C \\
        B & D & E \\
        C & E & F
    \end{bmatrix} > 0
\]

Finally, our conditions for a minimum are:
\begin{itemize}
    \item $$f_{xx} > 0$$.
    \item \[
        \det \begin{bmatrix}
            f_{xx} & f_{xy} \\
            f_{yx} & f_{yy}
        \end{bmatrix} > 0
    \]
    \item \[
        \det \begin{bmatrix}
            f_{xx} & f_{xy} & f_{xz} \\
            f_{yx} & f_{yy} & f_{yz} \\
            f_{zx} & f_{zy} & f_{zz}
        \end{bmatrix} > 0
    \]
\end{itemize}

And for a maximum:
\begin{itemize}
    \item $$f_{xx} < 0$$.
    \item \[
        \det \begin{bmatrix}
            f_{xx} & f_{xy} \\
            f_{yx} & f_{yy}
        \end{bmatrix} > 0
    \]
    \item \[
        \det \begin{bmatrix}
            f_{xx} & f_{xy} & f_{xz} \\
            f_{yx} & f_{yy} & f_{yz} \\
            f_{zx} & f_{zy} & f_{zz}
        \end{bmatrix} < 0
    \]
\end{itemize}

More generally, for a minimum we need these matrices, called the leading minors of the Hessian, to be positive and for a maximum we need them to alternate in sign.

\section{Lagrange Multipliers}
Suppose we wanted to minimise a function $f(x, y) = x^2 + y^2$ subject to a constraint, i.e. saying $x + y = 1$ (and therefore excluding the $(0, 0)$ trivial case).

While we could do this by writing $y = 1-x$, substituting and minimising solely in terms of $x$, there is a more sophisticated method that becomes useful in more complex cases where the constraint is not invertable.

We consider the function:
\[
    F(x, y, \lambda) = x^2 + y^2 + \lambda(x + y - 1)
\]

And minimise with respect to all three variables. When this happens, $F(x, y,\lambda) = f(x, y)$ at the minimised point.

Generally, to find the extrema of $f(x_1, x_2, \ldots, x_n)$ subject to $p$ constraints $g_i(x_1, x_2, \ldots, x_n) = 0$, where $p<n$, we introduce a new function of $n+p$ variables:

\[
    F(x_1, x_2, \ldots, x_n, \lambda_1, \lambda_2, \ldots, \lambda_p) = f(x_1, x_2, \ldots, x_n) + \sum_{i=1}^{p} \lambda_i g_i(x_1, x_2, \ldots, x_p)
\]

We then extremise $F$ wrt all $n+p$ variables. Extremising with respect to $\lambda_i$ enforces the constraint:
\[
    F_{\lambda_i} = g_i(x_1, x_2, \ldots, x_n) = 0
\]

So the the lambda term ends up dropping out and $F$ and $f$ have the same numerical value at extrema.